% p328_1; source PDF page 328; manually verified editable table.
\begin{tikzpicture}[x=1cm,y=-1cm,inner sep=0pt,outer sep=0pt]
\path[use as bounding box] (0,0) rectangle (16.000,7.560);
\draw[line width=.4pt] (0.000,0.000) rectangle (2.000,0.915);
\node[align=center,text width=1.740cm,font={\linespread{1}\fontsize{10}{13.50}\selectfont\bfseries}] at (1.000,0.458) {主观心理};
\draw[line width=.4pt] (2.000,0.000) rectangle (9.200,0.915);
\node[align=center,text width=6.940cm,font={\linespread{1}\fontsize{10}{13.50}\selectfont\bfseries}] at (5.600,0.458) {罪名};
\draw[line width=.4pt] (9.200,0.000) rectangle (16.000,0.915);
\node[align=center,text width=6.540cm,font={\linespread{1}\fontsize{10}{13.50}\selectfont\bfseries}] at (12.600,0.458) {范例};
\draw[line width=.4pt] (0.000,0.915) rectangle (2.000,5.445);
\node[align=center,text width=1.740cm,font={\linespread{1}\fontsize{10}{13.50}\selectfont}] at (1.000,3.180) {故意犯罪};
\draw[line width=.4pt] (2.000,0.915) rectangle (9.200,2.530);
\node[align=left,text width=6.940cm,font={\linespread{1}\fontsize{10}{13.50}\selectfont}] at (5.600,1.723) {危险达到放火罪、爆炸罪的相当程度：以危险方法危害公共安全罪};
\draw[line width=.4pt] (9.200,0.915) rectangle (16.000,2.530);
\node[align=left,text width=6.540cm,font={\linespread{1}\fontsize{10}{13.50}\selectfont}] at (12.600,1.723) {甲在高速公路上故意逆向高速行驶。甲构成以危险方法危害公共安全罪};
\draw[line width=.4pt] (2.000,2.530) rectangle (9.200,5.445);
\node[align=left,text width=6.940cm,font={\linespread{1}\fontsize{10}{13.50}\selectfont}] at (5.600,3.988) {危险未达到放火罪、爆炸罪的相当程度：危险驾驶罪、妨害安全驾驶罪};
\draw[line width=.4pt] (9.200,2.530) rectangle (16.000,5.445);
\node[align=left,text width=6.540cm,font={\linespread{1}\fontsize{10}{13.50}\selectfont}] at (12.600,3.988) {例1，甲在道路上醉酒驾驶，构成危险驾驶罪。\\{}例2，乘客短时间内抢夺公交车的方向盘，构成妨害安全驾驶罪};
\draw[line width=.4pt] (0.000,5.445) rectangle (2.000,7.560);
\node[align=center,text width=1.740cm,font={\linespread{1}\fontsize{10}{13.50}\selectfont}] at (1.000,6.503) {过失犯罪};
\draw[line width=.4pt] (2.000,5.445) rectangle (9.200,7.560);
\node[align=left,text width=6.940cm,font={\linespread{1}\fontsize{10}{13.50}\selectfont}] at (5.600,6.503) {具体罪名：交通肇事罪。\\{}兜底罪名（补充罪名）：过失以危险方法危害公共安全罪};
\draw[line width=.4pt] (9.200,5.445) rectangle (16.000,7.560);
\node[align=left,text width=6.540cm,font={\linespread{1}\fontsize{10}{13.50}\selectfont}] at (12.600,6.503) {甲驾车闯红灯，不慎撞死人，构成交通肇事罪};
\end{tikzpicture}
